\noindent
Let $R$ be a ring and let $A$, $B$ be $R$-algebras. In this section we
discuss $L_{A \otimes_R B/R}$. Most of the information we want is contained
in the following diagram
\begin{equation}
\label{equation-tensor-product}
\vcenter{
\xymatrix{
L_{A/R} \otimes_A^\mathbf{L} (A \otimes_R B) \ar[r] &
L_{A \otimes_R B/B} \ar[r] &
E \\
L_{A/R} \otimes_A^\mathbf{L} (A \otimes_R B) \ar[r] \ar@{=}[u] &
L_{A \otimes_R B/R} \ar[r] \ar[u] &
L_{A \otimes_R B/A} \ar[u] \\
 &
L_{B/R} \otimes_B^\mathbf{L} (A \otimes_R B) \ar[u] \ar@{=}[r] &
L_{B/R} \otimes_B^\mathbf{L} (A \otimes_R B) \ar[u]
}
}
\end{equation}
Explanation: The middle row is the fundamental triangle
(\ref{equation-triangle}) for the ring maps $R \to A \to A \otimes_R B$.
The middle column is the fundamental triangle
(\ref{equation-triangle}) for the ring maps $R \to B \to A \otimes_R B$.
Next, $E$ is an object of $D(A \otimes_R B)$ which ``fits'' into the
upper right corner, i.e., which turns both the top row
and the right column into distinguished triangles. Such an $E$
exists by Derived Categories, Proposition \ref{derived-proposition-9}
applied to the lower left square (with $0$ placed in the missing
spot). To be more explicit, we could for example define $E$ as the cone
(Derived Categories, Definition \ref{derived-definition-cone})
of the map of complexes
$$
L_{A/R} \otimes_A^\mathbf{L} (A \otimes_R B) \oplus
L_{B/R} \otimes_B^\mathbf{L} (A \otimes_R B)
\longrightarrow
L_{A \otimes_R B/R}
$$
and get the two maps with target $E$ by an application of TR3.
In the Tor independent case the object $E$ is zero.